\documentclass[11pt]{article}
\usepackage[margin=1in]{geometry}
\usepackage{xcolor}
\usepackage[breakable,listings]{tcolorbox}
\tcbuselibrary{listings}
\usepackage{hyperref}
\usepackage{enumitem}
\usepackage{listings}
\usepackage{amsmath}
\usepackage{amssymb}
\usepackage{booktabs}
\usepackage{array}
\usepackage{tabularx}
\usepackage{longtable}

% Define colors
\definecolor{osaorange}{RGB}{220,115,38}
\definecolor{aiblue}{RGB}{30,144,255}

% Configure hyperref
\hypersetup{
    colorlinks=true,
    linkcolor=blue,
    urlcolor=blue
}

% Code environment
\newtcblisting{rcode}{
  colback=white,
  colframe=gray!50,
  boxrule=0.5pt,
  arc=0pt,
  left=3pt,
  right=3pt,
  top=3pt,
  bottom=3pt,
  width=\textwidth,
  breakable,
  listing only,
  listing options={
    basicstyle=\small\ttfamily,
    breaklines=true,
    breakatwhitespace=true,
    columns=flexible,
    keepspaces=true,
    showstringspaces=false
  }
}

\title{}
\author{}
\date{}

\begin{document}

% Header box
\begin{tcolorbox}[colback=osaorange,colframe=black,boxrule=2pt,arc=0pt,left=3pt,right=3pt,top=6pt,bottom=6pt,boxsep=2pt]
\centering
\textcolor{white}{\textbf{\Large H524 Introduction to Biostatistics}}\\
\textcolor{white}{\textbf{\Large Final Group Project Specification}}\\
\textcolor{white}{\textbf{\Large Fall 2025}}
\end{tcolorbox}

\vspace{8pt}

\noindent\textbf{Project Timeline}\\
\textbf{Week 7-8:} Project preparation homework (dataset selection, proposal)\\
\textbf{Weeks 9-10:} Data analysis and report writing\\
\textbf{Week 10:} In-class presentations (10-15 minutes per group)\\
\textbf{End of Term:} Final written report and code submission\\
\textbf{Weight:} 25\% of final course grade

\section{Project Overview}

\subsection{Learning Objectives}

By completing this group project, you will:

\begin{itemize}
\item Apply statistical methods learned throughout the course to a real-world health dataset
\item Formulate and test meaningful research questions using appropriate statistical techniques
\item Create publication-quality tables, figures, and reports
\item Collaborate effectively in a team environment
\item Communicate statistical findings to both technical and non-technical audiences
\item \textcolor{aiblue}{\textbf{Integrate AI tools responsibly} to enhance your analysis while maintaining statistical rigor}
\item Develop reproducible research skills using R and RMarkdown
\end{itemize}

\subsection{Project Scope}

This is a \textbf{GROUP project} with 4-5 students per team. Each group will:

\begin{enumerate}
\item Select a publicly available health-related dataset (approved in Week 7-8)
\item Develop 2-3 research questions that can be addressed using methods from this course
\item Conduct a complete statistical analysis including descriptive statistics, hypothesis testing, and regression
\item Produce a professional written report (10-15 pages)
\item Deliver an in-class presentation (10-15 minutes)
\item Submit well-documented, reproducible R code
\end{enumerate}

\section{Required Deliverables}

\subsection{Deliverable 1: Written Report (10-15 pages)}

Your report should follow the structure of a scientific research article. Use RMarkdown to create a professional PDF document.

\subsubsection{Executive Summary (1 page)}

\begin{itemize}
\item Brief background and context
\item Key research questions
\item Main findings (2-3 bullet points)
\item Primary conclusions and implications
\end{itemize}

\subsubsection{Introduction \& Background (2 pages)}

\begin{itemize}
\item \textbf{Research Context:} Why is this topic important for public health?
\item \textbf{Literature Review:} Summarize 3-5 relevant published studies (use PubMed, Google Scholar)
\item \textbf{Research Questions:} State 2-3 specific, testable research questions
\item \textbf{Hypotheses:} For each research question, state your null and alternative hypotheses
\item \textbf{References:} Use APA format for all citations
\end{itemize}

\subsubsection{Data \& Methods (2-3 pages)}

\begin{itemize}
\item \textbf{Dataset Description:}
  \begin{itemize}
  \item Source and citation
  \item Study design (cross-sectional, cohort, etc.)
  \item Sample size and time period
  \item Population characteristics
  \end{itemize}

\item \textbf{Variables and Measures:}
  \begin{itemize}
  \item Primary outcome variable(s)
  \item Primary exposure/predictor variable(s)
  \item Confounding or effect modification variables
  \item Variable types (continuous, categorical, ordinal)
  \item Any transformations or recoding performed
  \end{itemize}

\item \textbf{Statistical Methods:}
  \begin{itemize}
  \item Descriptive statistics used
  \item Hypothesis tests performed (with justification)
  \item Regression models fitted
  \item Significance level ($\alpha$) used
  \item How missing data was handled
  \end{itemize}

\item \textbf{Software:}
  \begin{itemize}
  \item R version and key packages used
  \item \textcolor{aiblue}{AI tools used and their role in the analysis}
  \end{itemize}
\end{itemize}

\subsubsection{Results (3-4 pages)}

\begin{itemize}
\item \textbf{Descriptive Statistics:}
  \begin{itemize}
  \item Table 1: Participant characteristics (demographics, key variables)
  \item Summary statistics: means, SDs, medians, IQRs, frequencies, percentages
  \item Visual summaries: histograms, boxplots, bar charts
  \end{itemize}

\item \textbf{Primary Analyses:}
  \begin{itemize}
  \item Present results for each research question
  \item Include test statistics, p-values, and confidence intervals
  \item Provide effect sizes where appropriate
  \item Use tables and figures to communicate findings clearly
  \end{itemize}

\item \textbf{Table and Figure Requirements:}
  \begin{itemize}
  \item All tables numbered (Table 1, Table 2, etc.) with descriptive captions
  \item All figures numbered (Figure 1, Figure 2, etc.) with descriptive captions
  \item Tables: Use publication-quality formatting (no vertical lines, minimal horizontal lines)
  \item Figures: Clear axis labels, legends, and titles; readable font sizes
  \item Reference all tables/figures in the text
  \end{itemize}

\item \textbf{Statistical Reporting Standards:}
  \begin{itemize}
  \item Report exact p-values (not just ``p < 0.05'') when p > 0.001
  \item Include 95\% confidence intervals for key estimates
  \item Report effect sizes (e.g., mean differences, odds ratios, correlation coefficients)
  \end{itemize}
\end{itemize}

\subsubsection{Discussion \& Conclusions (2-3 pages)}

\begin{itemize}
\item \textbf{Interpretation of Findings:}
  \begin{itemize}
  \item What do your results mean in the context of your research questions?
  \item How do your findings compare to the literature?
  \item What are the public health implications?
  \end{itemize}

\item \textbf{Limitations:}
  \begin{itemize}
  \item Study design limitations
  \item Data quality issues
  \item Statistical limitations
  \item Generalizability concerns
  \end{itemize}

\item \textbf{Strengths:}
  \begin{itemize}
  \item What makes this analysis valuable?
  \item What did you do particularly well?
  \end{itemize}

\item \textbf{Future Directions:}
  \begin{itemize}
  \item What additional analyses could be done?
  \item What data would help answer remaining questions?
  \item What interventions or policies might be informed by this work?
  \end{itemize}

\item \textbf{Conclusions:}
  \begin{itemize}
  \item Concise summary of key findings
  \item Main take-home messages
  \end{itemize}
\end{itemize}

\subsubsection{References}

\begin{itemize}
\item Minimum of 5 references in APA format
\item Include dataset citation
\item Use primary sources (journal articles, government reports)
\end{itemize}

\subsubsection{Appendix (if needed)}

\begin{itemize}
\item Supplementary tables or figures
\item Additional sensitivity analyses
\item Extended technical details
\end{itemize}

\subsection{Deliverable 2: R Code (Well-Documented)}

Submit your complete R code as either:
\begin{itemize}
\item An RMarkdown (.Rmd) file that generates your report, OR
\item A standalone R script (.R) file with all analyses
\end{itemize}

\subsubsection{Code Requirements}

\begin{itemize}
\item \textbf{Reproducibility:}
  \begin{itemize}
  \item Code runs from start to finish without errors
  \item Includes code to load the dataset (with instructions on where to obtain it)
  \item Sets random seed if any randomization is used
  \item Includes \texttt{sessionInfo()} at the end
  \end{itemize}

\item \textbf{Organization:}
  \begin{itemize}
  \item Clear sections with headers (e.g., ``Data Loading'', ``Cleaning'', ``Analysis'', ``Figures'')
  \item Logical flow from data import to final results
  \item All packages loaded at the beginning
  \end{itemize}

\item \textbf{Documentation:}
  \begin{itemize}
  \item Comments explaining what each section does
  \item Comments for any complex or non-obvious code
  \item Variable names that are clear and descriptive
  \end{itemize}

\item \textbf{Good Coding Practices:}
  \begin{itemize}
  \item Consistent indentation and spacing
  \item No hard-coded values that should be variables
  \item Functions used for repeated operations
  \item Avoid unnecessary code repetition
  \end{itemize}

\item \textbf{AI-Generated Code:}
  \begin{itemize}
  \item \textcolor{aiblue}{Mark which sections were AI-generated (use comments)}
  \item \textcolor{aiblue}{Document how you verified AI code works correctly}
  \item \textcolor{aiblue}{Explain any modifications you made to AI output}
  \end{itemize}
\end{itemize}

\subsection{Deliverable 3: Presentation (10-15 minutes)}

Each group will present their project during Week 10.

\subsubsection{Presentation Structure}

\begin{itemize}
\item \textbf{Introduction (2-3 minutes):}
  \begin{itemize}
  \item Research question and motivation
  \item Brief background
  \end{itemize}

\item \textbf{Methods (2-3 minutes):}
  \begin{itemize}
  \item Dataset overview
  \item Key variables
  \item Statistical methods used
  \end{itemize}

\item \textbf{Results (4-6 minutes):}
  \begin{itemize}
  \item 2-3 key tables or figures
  \item Main statistical findings
  \item Focus on most important results
  \end{itemize}

\item \textbf{Discussion \& Conclusions (2-3 minutes):}
  \begin{itemize}
  \item What do findings mean?
  \item Limitations
  \item Take-home message
  \end{itemize}

\item \textbf{Q\&A (5 minutes):}
  \begin{itemize}
  \item Be prepared to answer questions about methods, results, and interpretations
  \end{itemize}
\end{itemize}

\subsubsection{Presentation Requirements}

\begin{itemize}
\item Create slides using Beamer (LaTeX) or PowerPoint
\item Each team member must present at least one section
\item Include a title slide with all team member names
\item Use clear, readable fonts (minimum 20pt for body text)
\item Limit text on slides (use bullet points, not paragraphs)
\item Include visual elements (tables, figures, diagrams)
\item Practice timing to stay within 10-15 minutes
\item Upload slides to Canvas before your presentation
\end{itemize}

\section{Statistical Analysis Requirements}

Your project must include \textbf{at least 4} of the following statistical methods:

\begin{enumerate}
\item \textbf{Descriptive Statistics and Visualizations}
  \begin{itemize}
  \item Summary statistics for all key variables
  \item At least 3 high-quality figures (histograms, boxplots, scatterplots, etc.)
  \item At least 1 well-formatted table of descriptive statistics
  \end{itemize}

\item \textbf{Confidence Intervals}
  \begin{itemize}
  \item 95\% CIs for means, proportions, or regression coefficients
  \item Proper interpretation of confidence intervals
  \end{itemize}

\item \textbf{Hypothesis Testing}
  \begin{itemize}
  \item Choose appropriate test(s): two-sample t-test, paired t-test, one-way ANOVA, chi-square test
  \item State null and alternative hypotheses
  \item Report test statistics, p-values, and conclusions
  \end{itemize}

\item \textbf{Correlation Analysis}
  \begin{itemize}
  \item Pearson or Spearman correlation
  \item Correlation matrix if multiple variables
  \item Scatterplots with correlation coefficients
  \end{itemize}

\item \textbf{Simple Linear Regression}
  \begin{itemize}
  \item Fit at least one regression model
  \item Interpret slope and intercept
  \item Report R-squared
  \item Check model assumptions (residual plots)
  \item Predict values for new observations
  \end{itemize}

\item \textbf{Additional Methods (as appropriate)}
  \begin{itemize}
  \item Stratified analyses
  \item Effect modification or interaction assessment
  \item Sensitivity analyses
  \item Non-parametric methods (if assumptions violated)
  \end{itemize}
\end{enumerate}

\textbf{Note:} Choose methods appropriate for your research questions and data. Quality of analysis is more important than quantity.

\section{AI Integration Requirements}

This course embraces AI as a tool for learning and productivity. Your project must document AI use.

\subsection{Required AI Documentation}

In a dedicated section of your report or in code comments, address:

\begin{enumerate}
\item \textbf{AI Tools Used:}
  \begin{itemize}
  \item Which AI tools did you use? (e.g., Microsoft Copilot, GitHub Copilot, ChatGPT)
  \item For what tasks? (e.g., code generation, debugging, explanation of concepts, literature search)
  \end{itemize}

\item \textbf{How AI Assisted Your Work:}
  \begin{itemize}
  \item Provide specific examples of AI prompts and outputs
  \item Explain how AI helped you solve problems or understand concepts
  \item What percentage of your code was AI-generated vs. written by you?
  \end{itemize}

\item \textbf{Verification of AI Outputs:}
  \begin{itemize}
  \item How did you verify AI-generated code was correct?
  \item Did you find any errors in AI outputs? How did you fix them?
  \item What manual checks did you perform?
  \end{itemize}

\item \textbf{Reflection on AI's Role:}
  \begin{itemize}
  \item What did AI do well?
  \item What did AI struggle with?
  \item How did using AI affect your learning?
  \item Would you use AI differently next time?
  \end{itemize}
\end{enumerate}

\subsection{AI Best Practices}

\begin{tcolorbox}[colback=aiblue!10,colframe=aiblue,title=AI Use Guidelines]

\textbf{Encouraged:}
\begin{itemize}
\item[+] Using AI to generate starter code templates
\item[+] Asking AI to explain statistical concepts or error messages
\item[+] Using AI to help debug code
\item[+] Asking AI to suggest visualization options
\item[+] Using AI to improve code documentation
\end{itemize}

\textbf{Required:}
\begin{itemize}
\item[\checkmark] Verify all AI outputs are correct
\item[\checkmark] Understand every line of code you submit
\item[\checkmark] Document AI use transparently
\item[\checkmark] Test AI-generated code thoroughly
\end{itemize}

\textbf{Prohibited:}
\begin{itemize}
\item[-] Submitting AI-generated work you don't understand
\item[-] Using AI to fabricate data or results
\item[-] Copying AI output without verification
\item[-] Failing to disclose AI use when required
\end{itemize}

\end{tcolorbox}

\section{Grading Rubric}

Total: 100 points = 25\% of course grade

\subsection{Written Report (60 points)}

\begin{table}[h]
\centering
\begin{tabular}{p{5cm}cp{7cm}}
\toprule
\textbf{Component} & \textbf{Points} & \textbf{Criteria} \\
\midrule
Introduction \& Background & 10 & Clear research questions, relevant literature review, well-motivated hypotheses \\
\midrule
Methods Description & 10 & Complete description of data, variables, and statistical methods; appropriate method selection \\
\midrule
Results Presentation & 15 & Clear tables and figures, appropriate statistical reporting, comprehensive results \\
\midrule
Statistical Correctness & 15 & Correct application of methods, accurate calculations, proper interpretation \\
\midrule
Discussion \& Conclusions & 10 & Thoughtful interpretation, acknowledgment of limitations, clear conclusions \\
\bottomrule
\end{tabular}
\end{table}

\subsection{R Code (20 points)}

\begin{table}[h]
\centering
\begin{tabular}{p{5cm}cp{7cm}}
\toprule
\textbf{Component} & \textbf{Points} & \textbf{Criteria} \\
\midrule
Reproducibility & 10 & Code runs without errors, produces all reported results, includes session info \\
\midrule
Code Quality \& Documentation & 10 & Well-organized, clearly commented, follows good practices, AI use documented \\
\bottomrule
\end{tabular}
\end{table}

\subsection{Presentation (15 points)}

\begin{table}[h]
\centering
\begin{tabular}{p{5cm}cp{7cm}}
\toprule
\textbf{Component} & \textbf{Points} & \textbf{Criteria} \\
\midrule
Content \& Clarity & 10 & Clear communication, effective visuals, appropriate content selection, timing \\
\midrule
Team Participation & 5 & All members participate, balanced contributions, ability to answer questions \\
\bottomrule
\end{tabular}
\end{table}

\subsection{Peer Evaluation (5 points)}

\begin{table}[h]
\centering
\begin{tabular}{p{5cm}cp{7cm}}
\toprule
\textbf{Component} & \textbf{Points} & \textbf{Criteria} \\
\midrule
Individual Contribution & 5 & Based on peer evaluations of contribution to team effort \\
\bottomrule
\end{tabular}
\end{table}

\textbf{Note:} Peer evaluation forms will be completed confidentially by each team member after the project is submitted.

\section{Group Dynamics}

\subsection{Team Formation}

\begin{itemize}
\item Groups of 4-5 students will be formed in Week 7
\item You may self-select groups or be assigned by instructor
\item All group members are expected to contribute equally
\end{itemize}

\subsection{Equal Contribution Expectation}

Each team member should contribute to:
\begin{itemize}
\item Dataset selection and proposal development
\item Literature review
\item Data analysis and coding
\item Writing different sections of the report
\item Creating presentation slides
\item Presenting during the oral presentation
\end{itemize}

\subsection{Peer Evaluation Process}

After project submission, each team member will complete a confidential peer evaluation form rating:
\begin{itemize}
\item Quality of contributions
\item Reliability and timeliness
\item Communication and collaboration
\item Overall effort level
\end{itemize}

Individual grades may be adjusted up or down based on peer evaluations if there is evidence of unequal contribution.

\subsection{Conflict Resolution}

If conflicts arise within your group:

\begin{enumerate}
\item \textbf{Team Discussion:} Try to resolve issues through open communication within the group
\item \textbf{Document Issues:} Keep records of team meetings, task assignments, and communications
\item \textbf{Seek Help Early:} Contact the instructor if issues cannot be resolved internally
\item \textbf{Formal Complaint:} If needed, submit written documentation of the issue to the instructor
\end{enumerate}

Waiting until the last minute to report conflicts is not acceptable. Notify the instructor as soon as problems arise.

\section{Technical Specifications}

\subsection{File Formats}

\begin{itemize}
\item \textbf{Written Report:} PDF (generated from RMarkdown preferred)
\item \textbf{Code:} .R or .Rmd files
\item \textbf{Presentation:} PDF or PowerPoint (.pptx)
\item \textbf{Data:} Include dataset or clear instructions for obtaining it
\end{itemize}

\subsection{File Naming Conventions}

Use the following naming scheme:

\begin{itemize}
\item Report: \texttt{GroupX\_Final\_Report.pdf}
\item Code: \texttt{GroupX\_Analysis\_Code.R} or \texttt{GroupX\_Analysis\_Code.Rmd}
\item Presentation: \texttt{GroupX\_Presentation.pdf}
\item Data (if included): \texttt{GroupX\_Data.csv} (or other format)
\end{itemize}

Replace ``X'' with your group number (e.g., Group1, Group2, etc.)

\subsection{Submission Method}

\begin{itemize}
\item Create a single ZIP file containing all required files
\item Name the ZIP file: \texttt{GroupX\_Final\_Project.zip}
\item Submit via Canvas by the deadline
\item Only ONE team member needs to submit (include all names in submission)
\end{itemize}

\subsection{Late Policy}

\begin{itemize}
\item \textbf{Presentations:} Cannot be made up if missed (0 points for presentation component)
\item \textbf{Written Report:} 10\% penalty per day late (up to 3 days)
\item \textbf{After 3 days:} No credit given
\item \textbf{Extensions:} Must be requested at least 48 hours before deadline with valid reason
\end{itemize}

\section{Example Projects}

To help you understand expectations, here are examples of appropriate project topics and approaches:

\subsection{Example 1: NHANES Physical Activity and Hypertension}

\textbf{Dataset:} National Health and Nutrition Examination Survey (NHANES)

\textbf{Research Questions:}
\begin{itemize}
\item Is physical activity level associated with prevalence of hypertension in US adults?
\item Does this association differ by age group?
\end{itemize}

\textbf{Methods Used:}
\begin{itemize}
\item Descriptive statistics for demographics and key variables
\item Chi-square test for association between physical activity category and hypertension
\item Two-sample t-test comparing blood pressure between active and sedentary individuals
\item Logistic regression predicting hypertension from physical activity (controlling for age, BMI)
\item Stratified analysis by age group
\end{itemize}

\textbf{Key Findings (hypothetical):}
\begin{itemize}
\item 32\% of sedentary adults had hypertension vs. 18\% of physically active adults (p < 0.001)
\item After adjusting for age and BMI, physical activity remained protective (OR = 0.58, 95\% CI: 0.45-0.75)
\item Association strongest in adults over 50 years old
\end{itemize}

\subsection{Example 2: Framingham Heart Study and Cholesterol}

\textbf{Dataset:} Framingham Heart Study (teaching dataset)

\textbf{Research Questions:}
\begin{itemize}
\item What is the relationship between age and total cholesterol?
\item Does smoking status affect cholesterol levels?
\end{itemize}

\textbf{Methods Used:}
\begin{itemize}
\item Descriptive statistics and histograms for cholesterol distribution
\item Correlation analysis between age and cholesterol
\item Simple linear regression: cholesterol $\sim$ age
\item Two-sample t-test comparing cholesterol between smokers and non-smokers
\item Multiple linear regression: cholesterol $\sim$ age + smoking + BMI
\end{itemize}

\textbf{Key Findings (hypothetical):}
\begin{itemize}
\item Moderate positive correlation between age and cholesterol (r = 0.42, p < 0.001)
\item For each additional year of age, cholesterol increases by 0.85 mg/dL (95\% CI: 0.70-1.00)
\item Smokers had 8.2 mg/dL higher cholesterol than non-smokers (p = 0.03)
\end{itemize}

\subsection{Example 3: COVID-19 Vaccination and Demographics}

\textbf{Dataset:} CDC COVID-19 vaccination data (county-level)

\textbf{Research Questions:}
\begin{itemize}
\item Is county-level vaccination rate associated with median household income?
\item Do urban and rural counties differ in vaccination rates?
\end{itemize}

\textbf{Methods Used:}
\begin{itemize}
\item Descriptive statistics for vaccination rates across counties
\item Scatterplot and correlation between vaccination rate and income
\item Simple linear regression: vaccination rate $\sim$ median income
\item Two-sample t-test comparing urban vs. rural vaccination rates
\item ANOVA comparing vaccination rates across regions (Northeast, South, Midwest, West)
\end{itemize}

\textbf{Key Findings (hypothetical):}
\begin{itemize}
\item Strong positive correlation between county income and vaccination rate (r = 0.68, p < 0.001)
\item Each \$10,000 increase in median income associated with 5.2\% higher vaccination rate
\item Urban counties had 12\% higher vaccination rates than rural counties (p < 0.001)
\end{itemize}

\section{Timeline and Milestones}

\begin{table}[h]
\centering
\begin{tabular}{p{3cm}p{10cm}}
\toprule
\textbf{Week} & \textbf{Milestone} \\
\midrule
Week 7 & Form groups; begin exploring datasets \\
\midrule
Week 8 & Submit project proposal (dataset, research questions, planned methods) \\
\midrule
Week 9 & Data cleaning and preliminary analyses; draft Methods section \\
\midrule
Week 10 (early) & Complete analyses; draft Results section; create presentation slides \\
\midrule
Week 10 (end) & IN-CLASS PRESENTATIONS \\
\midrule
End of Term & Final written report and code due \\
\bottomrule
\end{tabular}
\end{table}

\subsection{Week 8 Proposal Requirements}

Submit a 1-2 page proposal including:
\begin{itemize}
\item Group members and contact information
\item Dataset selected (with citation/link)
\item 2-3 proposed research questions
\item List of variables to be analyzed
\item Planned statistical methods
\item Potential challenges or concerns
\end{itemize}

This proposal will receive feedback from the instructor to help guide your project.

\section{Resources}

\subsection{RMarkdown Report Template}

An RMarkdown template for your report is provided on Canvas:
\begin{itemize}
\item \texttt{H524\_Project\_Report\_Template.Rmd}
\end{itemize}

This template includes:
\begin{itemize}
\item YAML header for PDF generation
\item Section structure matching this specification
\item Example code for tables and figures
\item Citation examples
\end{itemize}

\subsection{Beamer Presentation Template}

A Beamer (LaTeX) template for your presentation is provided on Canvas:
\begin{itemize}
\item \texttt{H524\_Project\_Presentation\_Template.tex}
\end{itemize}

This template includes OSU branding and example slides.

\subsection{Peer Evaluation Form}

The peer evaluation form will be distributed after presentations and is due with the final report. It includes:
\begin{itemize}
\item Rating scales for each team member's contributions
\item Open-ended questions about team dynamics
\item Confidential submission to instructor only
\end{itemize}

\subsection{Dataset Resources}

Recommended sources for health datasets:
\begin{itemize}
\item \textbf{CDC Data:} \url{https://data.cdc.gov/}
\item \textbf{NHANES:} \url{https://wwwn.cdc.gov/nchs/nhanes/}
\item \textbf{Framingham Heart Study:} Teaching datasets available on Canvas
\item \textbf{Behavioral Risk Factor Surveillance System (BRFSS):} \url{https://www.cdc.gov/brfss/}
\item \textbf{UCI Machine Learning Repository:} \url{https://archive.ics.uci.edu/ml/}
\item \textbf{Kaggle Datasets:} \url{https://www.kaggle.com/datasets}
\item \textbf{Data.gov:} \url{https://data.gov/}
\end{itemize}

\subsection{Style Guide for Tables and Figures}

\textbf{Tables:}
\begin{itemize}
\item Use horizontal lines sparingly (top, bottom, below headers)
\item Avoid vertical lines
\item Left-align text, right-align numbers
\item Include units in column headers
\item Use footnotes for clarifications
\item Caption goes above the table
\end{itemize}

\textbf{Figures:}
\begin{itemize}
\item Use clear, readable fonts (minimum 10pt)
\item Label all axes with units
\item Include legends when multiple groups shown
\item Use colorblind-friendly palettes
\item Caption goes below the figure
\item Export at high resolution (300 dpi for publication)
\end{itemize}

\subsection{APA Citation Examples}

\textbf{Journal Article:}
\begin{quote}
Smith, J., \& Jones, M. (2023). The effect of exercise on blood pressure: A meta-analysis. \textit{American Journal of Public Health}, \textit{113}(4), 456-467. https://doi.org/10.2105/AJPH.2023.123456
\end{quote}

\textbf{Dataset:}
\begin{quote}
Centers for Disease Control and Prevention. (2022). \textit{National Health and Nutrition Examination Survey 2017-2020} [Dataset]. U.S. Department of Health and Human Services. https://wwwn.cdc.gov/nchs/nhanes/
\end{quote}

\textbf{Website:}
\begin{quote}
World Health Organization. (2023, June 15). \textit{Cardiovascular diseases (CVDs)}. https://www.who.int/news-room/fact-sheets/detail/cardiovascular-diseases-(cvds)
\end{quote}

\subsection{AI Assistance}

\textcolor{aiblue}{\textbf{AI Tools Can Help With:}}
\begin{itemize}
\item Generating code templates for common analyses
\item Debugging error messages
\item Creating ggplot2 visualization code
\item Formatting tables with \texttt{kable} or \texttt{gt}
\item Explaining statistical concepts
\item Suggesting appropriate statistical tests
\item Improving code documentation
\item Finding relevant literature
\end{itemize}

\textcolor{aiblue}{\textbf{Remember:}} Always verify AI outputs and understand every line of code you submit!

\section{Frequently Asked Questions}

\textbf{Q: Can we use data we collected ourselves?}\\
A: Yes, with instructor approval. You must have IRB approval if human subjects data was collected for research purposes.

\textbf{Q: How long should our presentation be?}\\
A: 10-15 minutes, followed by 5 minutes for questions. Practice your timing.

\textbf{Q: Can we do more than the minimum 4 statistical methods?}\\
A: Absolutely! Quality analysis is encouraged. Just ensure you do justice to each method.

\textbf{Q: What if we can't meet as a group?}\\
A: Use Zoom, Google Docs, and other collaboration tools. Document your group meetings and task assignments.

\textbf{Q: Can we change our research question after the proposal?}\\
A: Minor changes are fine. Major changes require instructor approval.

\textbf{Q: How much code documentation is enough?}\\
A: Every analysis section should have comments. Someone unfamiliar with your project should be able to follow your code.

\textbf{Q: What if a team member isn't contributing?}\\
A: Document the issue and contact the instructor immediately. Don't wait until the end.

\textbf{Q: Can we use methods not covered in class?}\\
A: Advanced methods are welcome if you can explain them correctly. Consult with the instructor first.

\textbf{Q: How do we cite AI tools?}\\
A: Include in your methods or acknowledgments: ``We used [Tool Name] [Version] to assist with [specific tasks]. All outputs were verified by the authors.''

\textbf{Q: Is there a page limit?}\\
A: The report should be 10-15 pages (excluding references and appendices). Quality over quantity.

\section{Academic Integrity}

All work submitted must represent your team's own understanding and effort.

\begin{tcolorbox}[colback=red!10,colframe=red!80,title=Academic Integrity Violations]

The following constitute academic dishonesty:
\begin{itemize}
\item Submitting AI-generated content without verification or understanding
\item Fabricating or falsifying data or results
\item Plagiarizing text from sources without proper citation
\item Copying code from other groups without attribution
\item Having someone outside your group complete substantial portions of the project
\item Submitting work completed by a previous year's group
\end{itemize}

\textbf{Consequences:} Academic integrity violations will result in a zero on the project and possible referral to the Office of Student Conduct.

\end{tcolorbox}

\section{Summary Checklist}

Before submitting, verify you have:

\begin{itemize}
\item[$\square$] Written report (10-15 pages) with all required sections
\item[$\square$] Executive summary, introduction, methods, results, discussion, conclusions, references
\item[$\square$] At least 4 statistical methods appropriately applied
\item[$\square$] Publication-quality tables and figures with captions
\item[$\square$] R code that runs without errors and is well-documented
\item[$\square$] AI use documented in report and code
\item[$\square$] Presentation slides uploaded to Canvas
\item[$\square$] All files named correctly and zipped together
\item[$\square$] Peer evaluation form completed
\item[$\square$] All team members' names on report and presentation
\item[$\square$] Submitted by deadline
\end{itemize}

\vspace{12pt}

\begin{center}
\Large\textbf{Good luck with your project!}
\end{center}

\vspace{8pt}

\noindent Questions? Contact the instructor during office hours or via Canvas messaging.

\end{document}
